\documentclass[10pt,a4paper]{article}
\usepackage{graphicx}
\usepackage{mathtools}
\usepackage{amssymb}
\usepackage{microtype}
\usepackage{algorithm}
\usepackage{algpseudocodex}
\usepackage[margin=1in]{geometry}
\title{\vspace*{-1cm}Heavy-Tailed GPM Summary}
\author{Yi Hao}
\begin{document}
	\maketitle
	\section{Introduction}
	
	Gradient Projection Memory (GPM) has emerged as a powerful framework for continual learning by protecting the orthogonal subspaces containing representations of past tasks, thereby preventing catastrophic forgetting. However, traditional GPM assumes a classical statistical regime where feature representations are delocalised and bounded by finite variance. In modern, highly optimised deep networks, the empirical spectral density of the weight matrices typically exhibits Heavy-Tailed Self-Regularisation (HTSR). Under these heavy-tailed conditions ($\alpha < 2$), standard GPM becomes computationally inefficient and overly restrictive: a small number of massive outlier dimensions completely dominate the variance calculation, forcing GPM to allocate unnecessarily large, dense orthogonal subspaces that rapidly exhaust the network's capacity. To resolve this bottleneck, we introduce the Heavy-Tailed Gradient Projection Method (HTGPM), a spectral memory-management framework designed specifically for deep neural networks operating in the Heavy-Tailed Self-Regularisation (HTSR) regime. Rather than treating representations as flat geometric subspaces, HTGPM operates as a surgical spectral filter. By evaluating the system's Resolvent matrix just above the maximum eigenvalue, the algorithm leverages the principles of eigenstate localisation to isolate the exact structural hubs where task information concentrates. It then applies a soft power-law envelope to mask out background noise, compressing the allocated task bases and preserving maximum network plasticity for future tasks.
	
	\section{Methodology}
	
	The Heavy-Tailed Gradient Projection Method (HTGPM) processes task representations sequentially across layers. For a given layer, let $R \in \mathbb{R}^{N_{\text{dim}} \times M}$ represent the activation matrix generated by $M$ input images. The algorithm proceeds through four core mathematical phases:
	
	\subsection{1. Residual Space Extraction}
	Let $R \in \mathbb{R}^{N_{\text{dim}} \times M}$ be the activation matrix of a layer for the current task, where $N_{\text{dim}}$ is the neuron dimension and $M$ is the number of samples. If a basis from prior tasks $B_{\text{old}}$ exists, the incoming representations are projected onto its orthogonal complement to isolate novel task features:
	\begin{equation}
		R_{\text{proj}} = B_{\text{old}} (B_{\text{old}}^T R), \quad R_{\text{hat}} = R - R_{\text{proj}}
	\end{equation}
	We then compute the Singular Value Decomposition (SVD) of the residual space, $R_{\text{hat}} = U \Sigma V^T$. A base rank $k$ is determined by tracking the cumulative singular values required to meet a global variance energy threshold. This yields the raw, unweighted basis vectors $U_{\text{raw}} = U[:, :k]$.
	
	\subsection{2. Resolvent Construction and RMT Distance}
	Rather than relying on raw variance, HTGPM maps the collective spectral connectivity of the neurons. We form the correlation matrix $A = R_{\text{hat}} R_{\text{hat}}^T$ and evaluate its Resolvent matrix $G(z)$ (the Green's function) just beyond the spectral edge:
	\begin{equation}
		G(z) = (A - z I)^{-1}, \quad \text{where } z = \lambda_{\text{max}} + i\eta
	\end{equation}
	where $\lambda_{\text{max}}$ is the maximum eigenvalue of $A$, and $\eta$ is a small imaginary smoothing parameter. The envelope of the Resolvent is normalised into a bounded Random Matrix Theory (RMT) similarity metric, which defines our structural distance matrix:
	\begin{equation}
		D_{ij} = 1.0 - \frac{|G_{ij}|}{\sqrt{|G_{ii}| \cdot |G_{jj}|}}
	\end{equation}
	Two neurons structurally coupled via the collective spectrum yield a distance near $0$, while decoupled nodes approach $1$.
	
	\subsection{3. Dynamic Hub localisation via Inverse Participation Ratio}
	For each raw basis vector $v \in U_{\text{raw}}$, we compute its Inverse Participation Ratio (IPR) to determine its spatial concentration:
	\begin{equation}
		\text{IPR}(v) = \sum_{c=1}^{N_{\text{dim}}} v_c^4, \quad \text{yielding a dynamic hub allocation } k_{\text{hubs}} = \left\lceil \frac{1}{\text{IPR}(v)} \right\rceil
	\end{equation}
	The indices corresponding to the top $k_{\text{hubs}}$ largest squared amplitudes of $v$ are designated as feature hubs. For each hub index, we apply a heavy-tailed power-law localisation envelope across the RMT distance matrix to construct a soft spatial mask:
	\begin{equation}
		M_c = \max_{\text{hub}} \left( 1.0 + \frac{D_{\text{hub}, c}}{\alpha} \right)^{-\alpha}
	\end{equation}
	We apply this envelope as a multiplicative mask, $v_{\text{localised}} = \text{Normalise}(v \odot m)$, which suppresses neurons that are spectrally distant from the feature hubs. This effectively sparsifies the basis vector, focusing the feature representation onto the subset of neurons that share high-order correlations while pruning background noise.
	
	\subsection{4. Re-Orthogonalisation}
	Because the soft power-law masking breaks the strict orthogonality of the SVD, the stack of localised basis vectors is passed through a standard QR decomposition:
	\begin{equation}
		Q, R_{\text{qr}} = \text{QR}([v_{\text{localised}, 1}, \dots, v_{\text{localised}, k}])
	\end{equation}
	The signs are aligned using the diagonal of $R_{\text{qr}}$, and the resulting compact, orthogonal task basis $Q$ is concatenated to $B_{\text{old}}$ to preserve the memory footprint for subsequent tasks.
	
	\section{Discussion: Scaling Hypotheses and Future Steps}
	
	The theoretical foundation of HTGPM relies explicitly on the thermodynamic limit ($N \to \infty$) inherent to large random matrices. Consequently, we speculate that this algorithm exhibits a positive scaling law, where its performance advantage over standard GPM scales proportionally with network size. In low-dimensional regimes, such as small multi-layer perceptrons (MLPs), the algorithm is expected to underperform or fail; finite-size fluctuations destabilise the Resolvent operator's spectrum, and individual nodes are forced to be highly polyfunctional, preventing clear eigenstate localisation. Conversely, wide and deep architectures naturally undergo Anderson-like phase transitions, generating highly isolated structural hubs that our power-law envelope can cleanly exploit. 
	
	We explicitly acknowledge an algorithmic trade-off regarding computational complexity. Standard GPM requires only a fast, low-rank SVD at the end of a task. In contrast, HTGPM is noticeably more computationally expensive during task transitions due to the overhead of constructing the full covariance matrix, performing the matrix inversion for the complex Resolvent $G(z)$, and evaluating the power-law mask element-wise for each basis vector. However, this cost is strictly front-loaded at the boundaries between tasks; by isolating features into compact spectral corridors, HTGPM yields structurally compressed, sparse task bases. This drastically reduces the downstream memory footprint and prevents the rapid capacity saturation that typically plagues standard orthogonal projection methods over extended lifetimes.
	
	Future empirical validation will evaluate this framework on deep, wide architectures (e.g., Wide ResNets and Vision Transformers) initialised directly in a heavy-tailed state ($\alpha \approx 1.2$). We anticipate that optimising these networks will require a suppressed learning rate ($\eta_{\text{lr}}$) compared to standard Gaussian models to absorb the numerical shock of macro-step weight updates. However, we hypothesise that training velocity will remain high, as optimisation under heavy-tailed dynamics proceeds via non-local Lévy flights that tunnel directly between deep basins in the loss landscape. Our future tests will benchmark the sparsity and long-term memory capacity of HTGPM across extended task sequences to verify whether spectral localisation preserves network plasticity at scale.
	
	\newpage
	\section{Algorithm Implementation}
	\begin{algorithm}
		\caption{Heavy-Tailed Gradient Projection Method (HTGPM)}\label{alg:htgpm}
		\begin{algorithmic}[1]
			\Require Activation matrix $R$, Prior basis $B_{\text{old}}$, Tail index $\alpha$, Smoothing factor $\eta$
			\Ensure Updated memory basis $B_{\text{new}}$
			\Statex
			\State $R_{\text{hat}} \gets R - B_{\text{old}}(B_{\text{old}}^T R)$ \Comment{Extract novel residual features}
			\State $U, \Sigma, V^T \gets \text{SVD}(R_{\text{hat}})$ \Comment{Obtain raw base rank $k$ via energy threshold}
			\Statex 
			\State $G(z) \gets (R_{\text{hat}}R_{\text{hat}}^T - zI)^{-1} \text{ where } z = \lambda_{\text{max}} + i\eta$ \Comment{Evaluate spectral Resolvent at the edge}
			\State $D_{ij} \gets 1.0 - \frac{|G_{ij}|}{\sqrt{|G_{ii}| \cdot |G_{jj}|}}$ \Comment{Map Random Matrix Theory similarity metrics}
			\Statex
			\For{$\text{col\_idx} = 1$ to $k$}
			\State $v \gets U[:, \text{col\_idx}]$
			\State $k_{\text{hubs}} \gets \lceil 1 / \sum_{n} v_n^4 \rceil$ \Comment{Dynamic allocation via Inverse Participation Ratio}
			\State $\mathcal{H} \gets \text{argsort}(|v|^2)[\text{top } k_{\mathrm{hubs}}]$ \Comment{Isolate coordinate spatial hubs}
			\Statex
			\State $m \gets \max_{h \in \mathcal{H}} \left( 1.0 + D_{h, :} / \alpha \right)^{-\alpha}$ \Comment{Apply heavy-tailed localisation envelope}
			\State $v_{\text{sparse}} \gets \mathrm{Normalise}(v \odot m)$ \Comment{Suppress background spectral noise}
			\State Append $v_{\text{sparse}}$ to $W$
			\EndFor  % <--- THE CULPRIT: Fixed from \end{For} to \EndFor
			\Statex
			\State $Q, \sim \gets \text{QR}(W)$ \Comment{Restore strict orthogonality broken by masking}
			\State \Return $B_{\text{new}} \gets \mathrm{concatenate}([B_{\text{old}}, Q], \mathrm{axis}=1)$
		\end{algorithmic}
	\end{algorithm}
\end{document}